% !TEX root = ../main.tex
\section{Discussion}

\hypertarget{what-drives-the-inconsistency}{%
\subsection{What drives the inconsistency}\label{what-drives-the-inconsistency}}

The reproduction cleanly separates the ingredients of the phenomenon. The harvest-flow constraint is the operative between-period link: without it, the announced plan's tail remains optimal at each replan (the objective-gap diagnostic confirms the no-flow control is consistent at positive discount rates), and with it the re-solving planner strictly improves on -- or cannot even implement -- the announced tail. The non-declining-yield policy, the most restrictive, is the most frequently inconsistent. The disequilibrium forest structure is the fuel: a mature/old-growth landbase places a large pulse of immediately harvestable, financially over-mature volume in front of a regulated future that cannot sustain it, so the open-loop plan's promised smooth flow can only be constructed by commitments the future planner will not keep. The negatively valued strata are the tell: because CM-CE area enters the solution only to mitigate the binding flow constraint and is then abandoned under replanning, its presence in the open-loop solution is a direct, observable marker of the inconsistency rather than a real management intention.

\hypertarget{the-discount-rate-nuance}{%
\subsection{The discount-rate nuance}\label{the-discount-rate-nuance}}

One finding deserves emphasis because it is easy to misstate. The naive account -- ``a discounted objective plus an even-flow constraint breaks Bellman's principle'' -- suggests that the discount rate is the operative driver and that consistency could be restored by lowering it. The reproduction shows this is not so: genuine dynamic inconsistency -- the announced plan's tail being strictly suboptimal or infeasible from the realized state, per the objective-gap diagnostic (Section 3.3) -- occurs at every discount rate (0\%, 2\%, 4\%, and 6\%) for the flow-constrained policies. The resolution is in the thesis itself \citep[ch.~3, p.~58--59]{daugherty1991credibility}: because the discount factor is uniform across periods, it cancels out of the conditions that a plan must satisfy to be consistent, so the \emph{occurrence} of inconsistency does not require a particular rate. Discounting shapes \emph{which} plan is chosen and the \emph{magnitude} of the divergence, however: the mean relative divergence and total-volume shortfall are markedly larger at a 0\% rate than at 2--6\%. We note one subtlety the metric alone would obscure: at a 0\% rate the flat objective admits near-tie optima, so the trajectory-divergence metric also flags the no-flow control (whose announced tail the objective-gap diagnostic shows is largely still optimal there -- alternate optima, not genuine inconsistency); at positive rates the objective discriminates and the flow-constraint divergence is genuine. So the rate matters to the plan and to how badly the plan is wrong, but not to whether the plan is wrong. We present this as a reproduced nuance of the thesis, not as a contradiction of it -- recovering this counter-intuitive rate-independence on independent software is a meaningful check that the reproduction has captured the thesis's mechanism. The practical corollary is uncomfortable: discount-rate policy is not a remedy for non-credible plans.

\hypertarget{implications-for-the-literature}{%
\subsection{Implications for the literature}\label{implications-for-the-literature}}

The implication for the large literature of open-loop LP forest-planning studies -- and for planning practice built on single-solve strategic models -- is direct. Where a model combines a harvest-flow constraint with a disequilibrium initial forest, its projected harvest trajectories, residual growing-stock trajectories, and shadow prices are properties of a plan that sequential re-optimization will not follow, and trade-off analysis conducted on them is biased. The ``declining non-declining yield'' \citep[pp.~331--332]{gunn2007models} names the symptom exactly: the promised non-declining flow declines precisely because each successive planner declines to honour it. The same mechanism surfaces downstream, in the hierarchical link between strategic and operational planning: \citet{paradis2013risk} showed that incoherent linkage between long-term wood-supply allocation and short-term harvest planning induces a systematic drift between planned and implemented harvest (a divergence of forest system state), citing \citet{daugherty1991credibility} and naming the declining-non-declining-yield phenomenon. The dynamic inconsistency reproduced here is the analytic root of that practical drift, a connection to the longer economics literature on the even-flow constraint and the allowable-cut effect \citep{thompson1966traditional,schweitzer1972allowable,luckert1995allowable,hegan2000economic}. That the underlying analysis was never peer-reviewed and has been held only as an archival thesis plausibly explains why the phenomenon has been repeatedly rediscovered; an open, citable, reproducible reference implementation is useful to the field.

\hypertarget{remedies}{%
\subsection{Remedies}\label{remedies}}

The thesis and the subsequent literature point to three families of remedies, which the reproduction makes concrete. First, \emph{consistent formulations}: rather than solving one open-loop problem, compute solutions that are optimal for the sequential game the planning process actually is -- subgame-perfect or feedback solutions in which each period's decision is optimal given the state and the continuation policy. The thesis discusses the generation of consistent solutions; computing them exactly for the full case study is a natural follow-on enabled by this stack. Second, \emph{credible precommitment}: inconsistency arises because the future planner is free to re-solve; institutional commitments that genuinely bind future decisions -- for example, regeneration mandates attached to harvest -- change the future subproblem and can restore credibility to announced plans; the principal-agent literature on incentive-compatible timber-supply policy after disturbance develops this incentive-compatibility angle \citep{bogle2015protecting}. Third, \emph{receding-horizon practice with honest reporting}: if planning is in fact sequential replanning, then the simulator used here -- not the single open-loop solve -- is the honest model of what the planning process will deliver, and plan documents should report the realized replanned trajectory rather than the open-loop projection.

\hypertarget{limitations}{%
\subsection{Limitations}\label{limitations}}

Several limitations bound the claims. First, the data vintage: the case study is built from the 1990 FEIS (the FORPLAN model lineage with updated 1980s prices), cross-checked against the 1987 DEIS that Daugherty used directly; the two vintages' yield tables agree, and the thesis's exact base-case schedule is not matched period by period because the thesis's per-period realized schedule is not printed. Validation is against the thesis's printed anchors: the mature-type PNVs match Table 5.4 exactly (their volumes are derived from the anchors, cross-checked against the independent FEIS standing-volume curves), and the managed prescriptions match Table 5.3 in \emph{sign} and broad rotation range (the exact optimal-rotation ordering is only partially reproduced -- the FEIS prescription mapping simplifies some treatments). Second, the formulation caveat: the thesis used a Model II formulation; the reproduction uses Model I (Section 2.2), because the Model II flow-constraint compilation is not yet implemented in \texttt{ws3} -- a forced substitution whose behavioral equivalence under sequential replanning we note but do not demonstrate. Third, the thesis used ending-period (terminal) constraints with its flow-constrained runs to minimize horizon effects; we implemented these in the stack (currently opt-in and experimental) and found that with the corrected case-study data the model does not exhibit the horizon-end liquidation they guard against, so they are omitted from the reported grid. Fourth, the simulator is a rolling-horizon (receding-horizon) institution, which is how planning actually proceeds; a shrinking-horizon variant (the pure Bellman tail over a fixed terminal date) shows the same qualitative phenomenon at somewhat smaller magnitude (on landbase 1 under NDY at 4\%, mean relative divergence 0.050 shrinking vs 0.074 rolling; both inconsistent). Fifth, scope: this is one case study. The pervasiveness result (73\% of flow-constrained cells) is a property of this grid, and extension to other forests, formulations, and consistent-solution constructs is future work -- work that the open stack is intended to make easy.
